% ============================================================
% content/solutions.tex —— 2022 湖南卷物理 · 参考答案与详细解析
% 解析类型：source_embedded（转录自来源站 Word 版内嵌解析，独立复核）
% 复核记录见 metadata.yaml 的 notes
% ============================================================

\begin{answer}
\gksolhead{1}{单项选择题}
\gkanswer{答案 C}
卢瑟福的核式结构模型解释了 $\alpha$ 粒子散射实验现象，不能解释原子光谱的分立特征，原子光谱的分立特征由玻尔理论解释，故 A 错误；玻尔的原子理论只成功地解释了氢原子光谱，并未完全揭示微观粒子运动的规律，故 B 错误；光电效应表明光具有粒子性，故 C 正确；电子束穿过铝箔后的衍射图样说明电子具有波动性，故 D 错误。
\end{answer}

\begin{answer}
\gksolhead{2}{单项选择题}
\gkanswer{答案 A}
四根绝缘棒完全相同且对称放置，$a$ 与 $c$、$b$ 与 $d$ 分别关于几何中心 $O$ 对称，它们在 $O$ 点产生的电场两两等大反向，故四根棒都存在时 $O$ 点的电场强度为零。

移去 $a$ 处的绝缘棒后，$b$、$d$ 两棒在 $O$ 点产生的电场相互抵消，$O$ 点的电场等于 $c$ 棒单独产生的电场。$c$ 棒带正电，其电场方向沿 $O$ 与 $c$ 的连线背离 $c$，即垂直指向 $a$，故 A 正确、B 错误。

取无穷远处电势为零，四根带正电的棒在 $O$ 点产生的电势均为正值。移去带正电的 $a$ 棒后，$O$ 点处正电荷的电势能减少，由 $\varphi=E_{p}/q$ 可知 $O$ 点的电势减小，故 C、D 错误。
\end{answer}

\begin{answer}
\gksolhead{3}{单项选择题}
\gkanswer{答案 D}
导线静止时受重力 $G$、悬线拉力 $F_{T}$ 和安培力 $F_{\text{安}}$ 作用而平衡。磁场方向垂直指向 $OO'$，悬线也指向 $OO'$，故安培力 $F_{\text{安}}$ 与拉力 $F_{T}$ 垂直，三力构成直角三角形，$G$ 为斜边，如图所示。

\gkfig{figures/s3-force}

由几何关系得
\[
  F_{\text{安}}=G\sin\theta,\qquad F_{T}=G\cos\theta .
\]
又 $F_{\text{安}}=BIl$，故
\[
  \sin\theta=\frac{BIl}{G},
\]
即 $\sin\theta$ 与电流 $I$ 成正比，故 C 错误、D 正确。

电流 $I$ 增大时 $\theta$ 增大，$\cos\theta$ 减小，悬线拉力 $F_{T}$ 减小，由牛顿第三定律，导线对悬线的拉力也减小，故 B 错误。

当导线静止在题图甲右侧位置时，安培力必须水平向右才能与拉力的水平分力平衡。由左手定则可知，导线中电流方向由 $M$ 指向 $N$，故 A 错误。
\end{answer}

\begin{answer}
\gksolhead{4}{单项选择题}
\gkanswer{答案 B}
中子的质量与氢核质量相等，且小于氮核质量，即 $m_{\text{中}}=m_{\text{氢}}<m_{\text{氮}}$。

弹性正碰中，质量为 $m_{1}$ 的粒子以速度 $v_{0}$ 碰静止的质量为 $m_{2}$ 的粒子，碰后后者速度为
\[
  v_{2}'=\frac{2m_{1}}{m_{1}+m_{2}}v_{0}.
\]
中子与氢核质量相等，故碰后氢核速度
\[
  v_{1}=v_{0};
\]
中子与氮核相碰时
\[
  v_{2}=\frac{2m_{\text{中}}}{m_{\text{中}}+m_{\text{氮}}}v_{0}<v_{0},
\]
所以 $v_{1}>v_{2}$，C、D 均错误。

碰后氢核的动量
\[
  p_{\text{氢}}=m_{\text{氢}}v_{1}=m_{\text{中}}v_{0},
\]
碰后氮核的动量
\[
  p_{\text{氮}}=m_{\text{氮}}v_{2}=\frac{2m_{\text{氮}}m_{\text{中}}}{m_{\text{中}}+m_{\text{氮}}}v_{0}>m_{\text{中}}v_{0},
\]
故 $p_{\text{氮}}>p_{\text{氢}}$，A 错误。

碰后氢核的动能
\[
  E_{\text{k氢}}=\frac{1}{2}m_{\text{中}}v_{0}^{2},
\]
碰后氮核的动能
\[
  E_{\text{k氮}}=\frac{1}{2}m_{\text{氮}}v_{2}^{2}
  =\frac{2m_{\text{氮}}m_{\text{中}}^{2}}{(m_{\text{中}}+m_{\text{氮}})^{2}}v_{0}^{2}.
\]
两者之比
\[
  \frac{E_{\text{k氮}}}{E_{\text{k氢}}}
  =\frac{4m_{\text{中}}m_{\text{氮}}}{(m_{\text{中}}+m_{\text{氮}})^{2}}<1,
\]
故 $E_{\text{k氮}}<E_{\text{k氢}}$，B 正确。
\end{answer}

\begin{answer}
\gksolhead{5}{单项选择题}
\gkanswer{答案 A}
在飘带上任取一小段，设其长度为 $\Delta l$、质量为 $\Delta m$。由于飘带粗细一致、匀质，$\Delta m$ 与 $\Delta l$ 成正比，该小段所受重力 $\Delta m g$ 与长度成正比。风速水平向右、大小恒定且不随高度改变，飘带各处单位长度所受风力的大小和方向都相同，也与长度成正比。

因此，飘带上每一小段所受重力与风力的合力，其大小与长度成正比、方向处处相同，即飘带处在方向确定的等效均匀力场中。柔软绳在均匀力场中平衡时，必须沿合力方向伸直；否则合力会在垂直于绳的方向产生分力，使绳弯曲而不能平衡。所以稳定时飘带是一条倾斜的直线，最接近 A 图的形态。
\end{answer}

\begin{answer}
\gksolhead{6}{单项选择题}
\gkanswer{答案 B}
$\mathrm{P}_{1}$ 在副线圈正中间时，原、副线圈接入电路的匝数之比 $n_{1}:n_{2}=2:1$。设原线圈中电流为 $I$，由理想变压器电流关系 $n_{1}I=n_{2}I_{2}$ 得副线圈电流
\[
  I_{2}=2I .
\]
副线圈两端电压 $U_{2}=I_{2}R_{1}=2IR$，原线圈两端电压
\[
  U_{1}=\frac{n_{1}}{n_{2}}U_{2}=4IR .
\]
设输入电压有效值为 $U_{0}$，则
\[
  U_{0}=U_{1}+IR_{2}=4IR+IR_{2},
\]
解得
\[
  I=\frac{U_{0}}{4R+R_{2}} .
\]

保持 $\mathrm{P}_{1}$ 位置不变，$\mathrm{P}_{2}$ 向左滑动时 $R_{2}$ 减小，$I$ 增大，$U_{1}=4IR$ 增大，电压表示数 $U=U_{0}-U_{1}$ 减小（电压表测 $R_{2}$ 两端电压），故 A 错误。$R_{1}$ 消耗的功率
\[
  P_{1}=I_{2}^{2}R=4I^{2}R
\]
随 $I$ 增大而增大，故 B 正确。

保持 $\mathrm{P}_{2}$ 位置不变（$R_{2}=9R$），$\mathrm{P}_{1}$ 向下滑动时，副线圈接入匝数减少，$n=n_{1}/n_{2}$ 增大。同理
\[
  U_{2}=nIR,\qquad U_{1}=n^{2}IR,
\]
\[
  U_{0}=n^{2}IR+9IR
  \;\Longrightarrow\;
  I=\frac{U_{0}}{(n^{2}+9)R},
\]
$n$ 增大则 $I$ 减小，$U=9IR$ 减小，故 C 错误。

此时 $R_{1}$ 消耗的功率
\[
  P_{1}=(nI)^{2}R
  =\frac{n^{2}U_{0}^{2}R}{(n^{2}+9)^{2}R^{2}}
  =\frac{U_{0}^{2}}{R}\cdot\frac{n^{2}}{(n^{2}+9)^{2}} .
\]
令 $x=n^{2}$，则 $f(x)=\dfrac{x}{(x+9)^{2}}$，由
\[
  f'(x)=\frac{(x+9)^{2}-2x(x+9)}{(x+9)^{4}}
  =\frac{9-x}{(x+9)^{3}}
\]
知 $x=9$（即 $n=3$）时 $P_{1}$ 最大。初始 $n=2$，$\mathrm{P}_{1}$ 向下滑动时 $n$ 由 2 增大，$P_{1}$ 先增大后减小，故 D 错误。
\end{answer}

\begin{answer}
\gksolhead{7}{多项选择题}
\gkanswer{答案 AC}
由 $v$--$t$ 图像可知，在 $0\sim t_{1}$ 时间内返回舱速度减小，图线斜率的绝对值逐渐减小，加速度逐渐减小，故 B 错误。重力的瞬时功率
\[
  P=mgv
\]
随速度减小而减小，故 A 正确。

在 $t_{1}\sim t_{2}$ 时间内速度继续减小，质量不变，动量 $p=mv$ 随时间减小，故 C 正确。

在 $t_{2}\sim t_{3}$ 时间内速度不变，动能不变；返回舱高度下降，重力势能减小，故机械能减小，D 错误。
\end{answer}

\begin{answer}
\gksolhead{8}{多项选择题}
\gkanswer{答案 CD}
由开普勒第三定律
\[
  \frac{r_{\text{火}}^{3}}{T_{\text{火}}^{2}}=\frac{r_{\text{地}}^{3}}{T_{\text{地}}^{2}},
  \qquad r_{\text{火}}=\frac{3}{2}r_{\text{地}},
\]
得
\[
  T_{\text{火}}=\left(\frac{3}{2}\right)^{3/2}T_{\text{地}}
  =\sqrt{\frac{27}{8}}\,T_{\text{地}},
\]
故 A 错误（A 项给出的是其倒数）。

由万有引力提供向心力
\[
  G\frac{Mm}{r^{2}}=m\omega^{2}r
  \;\Longrightarrow\;
  \omega=\sqrt{\frac{GM}{r^{3}}},
\]
轨道半径越大角速度越小，故 $\omega_{\text{火}}<\omega_{\text{地}}$。

冲日时太阳、地球、火星共线，地球位于太阳与火星之间。地球公转角速度较大，从内侧“追上并超过”火星，在地球上观察，火星相对恒星背景向后退（由东向西），即逆行，故 B 错误、C 正确。

冲日时地球与火星的速度方向相同，相对速度大小
\[
  |\Delta v|=|v_{\text{地}}-v_{\text{火}}|,
\]
为两者速率之差，在公转一周中最小，故 D 正确。
\end{answer}

\begin{answer}
\gksolhead{9}{多项选择题}
\gkanswer{答案 BC}
发动机关闭、飞行器匀速下落时，阻力与重力平衡：
\[
  kv_{1}^{2}=mg,\qquad v_{1}=\SI{10}{\metre\per\second},
\]
得
\[
  k=\frac{mg}{100}.
\]
发动机以最大推力向上、飞行器匀速上升时
\[
  F_{\max}=mg+kv_{2}^{2}
  =mg+\frac{mg}{100}\times25
  =1.25mg,
\]
故 A 错误。

飞行器以 \SI{5}{\metre\per\second} 匀速水平飞行时，阻力 $kv_{3}^{2}=\dfrac{mg}{100}\times25=0.25mg$ 沿水平方向，重力 $mg$ 竖直向下，推力与二者的合力等大反向，受力分析如图所示。

\gkfig{figures/s9-force}

故
\[
  F=\sqrt{(mg)^{2}+(0.25mg)^{2}}
  =\frac{\sqrt{17}}{4}mg,
\]
B 正确。

发动机以最大推力推动飞行器匀速水平飞行时，设速率为 $v_{4}$，则
\[
  F_{\max}^{2}=(mg)^{2}+(kv_{4}^{2})^{2}.
\]
代入 $F_{\max}=1.25mg$ 得
\[
  kv_{4}^{2}=mg\sqrt{1.25^{2}-1}=0.75mg,
\]
又 $kv_{1}^{2}=mg$，故
\[
  v_{4}=v_{1}\sqrt{0.75}=5\sqrt{3}~\si{\metre\per\second},
\]
C 正确。

加速度最大时，最大推力竖直向下，飞行器以 \SI{5}{\metre\per\second} 的速率向上减速，阻力也竖直向下，由牛顿第二定律
\[
  F_{\max}+mg+kv_{5}^{2}=ma,
\]
\[
  a=\frac{1.25mg+mg+0.25mg}{m}=2.5g<3g,
\]
故 D 错误。
\end{answer}

\begin{answer}
\gksolhead{10}{多项选择题}
\gkanswer{答案 BD}
导体棒 $b$ 刚要滑动时，所受静摩擦力达到最大值，安培力
\[
  F_{\text{安}}=\mu mg=0.1\times0.1\times10~\si{\newton}=\SI{0.1}{\newton}.
\]
又 $F_{\text{安}}=BI_{b}l$，得
\[
  I_{b}=\frac{0.1}{0.1\times1}~\si{\ampere}=\SI{1}{\ampere}.
\]
导体棒 $b$ 与电阻 $R$ 并联且阻值均为 \SI{0.1}{\ohm}，故通过 $R$ 的电流 $I_{R}=I_{b}=\SI{1}{\ampere}$，干路电流
\[
  I_{\text{总}}=I_{b}+I_{R}=\SI{2}{\ampere}.
\]
电路总电阻
\[
  R_{\text{总}}=R_{a}+\frac{R_{b}R}{R_{b}+R}
  =0.1+\frac{0.1\times0.1}{0.1+0.1}~\si{\ohm}
  =\SI{0.15}{\ohm}.
\]
此时导体棒 $a$ 产生的感应电动势
\[
  E=I_{\text{总}}R_{\text{总}}=2\times0.15~\si{\volt}=\SI{0.3}{\volt},
\]
由 $E=Blv$ 得
\[
  v=\frac{E}{Bl}=\frac{0.3}{0.1\times1}~\si{\metre\per\second}=\SI{3}{\metre\per\second}.
\]

导体棒 $a$ 离开导轨后做平抛运动，下落时间
\[
  t=\sqrt{\frac{2h}{g}}=\sqrt{\frac{2\times0.8}{10}}~\si{\second}=\SI{0.4}{\second},
\]
水平位移
\[
  x=vt=3\times0.4~\si{\metre}=\SI{1.2}{\metre}\neq\SI{0.6}{\metre},
\]
故 A 错误。

离开导轨后导体棒 $a$ 的水平速度不变，且始终与磁场垂直，由 $E=Blv$ 知感应电动势不变，故 B 正确。

导体棒 $a$ 向右运动，由右手定则可判知电路中感应电流的方向，再由左手定则可知导体棒 $b$ 所受安培力的方向向左，$b$ 有向左运动的趋势，故 C 错误。

通过电路的总电荷量
\[
  q=\bar{I}\Delta t=\frac{\Delta\Phi}{R_{\text{总}}}
  =\frac{Bl\Delta x}{R_{\text{总}}}
  =\frac{0.1\times1\times1.74}{0.15}~\si{\coulomb}
  =\SI{1.16}{\coulomb}.
\]
$b$ 与 $R$ 并联且阻值相等，电流相等，故通过 $R$ 的电荷量
\[
  q_{R}=\frac{q}{2}=\SI{0.58}{\coulomb},
\]
D 正确。

\begin{gktip}
复核说明：独立复核时曾按楞次定律误判 $b$ 所受安培力向右、答案应为 BCD；经查新浪教育公布的 2022 湖南卷参考答案（7--10：AC、CD、BC、BD），第 10 题官方答案为 BD，已按官方答案更正。
\end{gktip}
\end{answer}

\begin{answer}
\gksolhead{11}{实验题}
\gkanswer{答案 (3) 见图线；(4) 15.35；(5) 127}
(3) 以 $n$ 为横轴、$l$ 为纵轴描点，五个点近似在同一直线上，用直线拟合，如图所示。

\gkfig{figures/s11-graph}

(4) 刻度尺的分度值为 \SI{1}{\milli\metre}，图丙中橡皮筋下端对应的读数为
\[
  l=15~\si{\centi\metre}+3.5~\si{\milli\metre}
   =153.5~\si{\milli\metre}=15.35~\si{\centi\metre}.
\]

(5) 设橡皮筋的劲度系数为 $k$、原长为 $x_{0}$。由胡克定律，挂 $n$ 枚硬币时
\[
  nm_{0}g=k(l-x_{0}),
\]
即 $l$ 与 $n$ 成线性关系，图线的斜率
\[
  \frac{\Delta l}{\Delta n}
  =\frac{16.56-10.51}{25-5}~\si{\centi\metre/\text{枚}}
  =0.3025~\si{\centi\metre/\text{枚}}.
\]
每增加一枚硬币，橡皮筋多伸长 $\Delta l=0.3025$~cm，故
\[
  k=\frac{m_{0}g}{\Delta l}
  =\frac{6.05\times10^{-3}g}{3.025\times10^{-3}}~\si{\newton\per\metre}
  =2.0g~\si{\newton\per\metre}.
\]
图线在纵轴上的截距
\[
  x_{0}=10.51-5\times0.3025~\si{\centi\metre}=9.00~\si{\centi\metre}.
\]
挂冰墩墩玩具时，同理
\[
  m_{1}g=k(l-x_{0}),
\]
\[
  m_{1}=\frac{k(l-x_{0})}{g}
  =\frac{2.0g\times(15.35-9.00)\times10^{-2}}{g}~\si{\kilo\gram}
  =0.127~\si{\kilo\gram}=127~\si{\gram}.
\]
\end{answer}

\begin{answer}
\gksolhead{12}{实验题}
\gkanswer{答案 (1) 大于；(2) $\times10$；(3) 向上；(4) 400}
(1) 由 $R_{m}<R_{n}$ 知，开关接 $m$ 时电路总电阻较小，电源电动势不变，由闭合电路欧姆定律 $I=E/R$ 知总电流较大，电流表 $\mathrm{G}$ 的示数 $I_{m}>I_{n}$，填“大于”。

(2) 开关接 $n$ 时欧姆表内阻较大。欧姆表内阻越大，同一指针偏转对应的待测电阻越大，倍率越大，故为“$\times10$”挡。

(3) 从“$\times1$”挡换成“$\times10$”挡，开关由 $m$ 拨向 $n$，电路总电阻增大，干路电流减小。①②短接调零时，需使电流表 $\mathrm{G}$ 满偏，即增大流过 $\mathrm{G}$ 所在支路的电流占比。$R_{0}$ 的滑片向上调节时，$R_{0\text{上}}$ 减小、$R_{0\text{下}}$ 增大，$\mathrm{G}$ 所在上支路电阻减小而下支路电阻增大，流过 $\mathrm{G}$ 的电流占比增大，故填“向上”。

(4) 设“$\times10$”挡时欧姆表内阻为 $R_{\text{内}}$，电流表满偏电流为 $I_{g}$，上支路电阻为 $R_{A}=R_{0\text{上}}+R_{G}$，下支路电阻为 $R_{B}=R_{0\text{下}}+R_{s}$。两支路并联，电压相等。电路可简化为图所示的结构：

\gkfig{figures/s12-circuit}

\begin{equation*}
  I_{G}R_{A}=I_{B}R_{B}
  \;\Longrightarrow\;
  I_{B}=I_{G}\frac{R_{A}}{R_{B}} .
\end{equation*}
于是干路电流
\[
  I_{\text{总}}=I_{G}+I_{B}=I_{G}\left(1+\frac{R_{A}}{R_{B}}\right).
\]
记 $K=I_{g}\left(1+\dfrac{R_{A}}{R_{B}}\right)$。①②短接时满偏：
\[
  \frac{E}{R_{\text{内}}}=K .
\]
接入 $R_{1}=\SI{100}{\ohm}$ 时指针为满偏的 $\dfrac{2}{3}$：
\[
  \frac{E}{R_{\text{内}}+100}=\frac{2}{3}K .
\]
接入 $R_{x}$ 时指针为满偏的 $\dfrac{1}{3}$：
\[
  \frac{E}{R_{\text{内}}+R_{x}}=\frac{1}{3}K .
\]
由前两式得
\[
  \frac{R_{\text{内}}+100}{R_{\text{内}}}=\frac{3}{2}
  \;\Longrightarrow\;
  R_{\text{内}}=\SI{200}{\ohm}.
\]
由第一、三两式得
\[
  \frac{R_{\text{内}}+R_{x}}{R_{\text{内}}}=3
  \;\Longrightarrow\;
  R_{x}=2R_{\text{内}}=\SI{400}{\ohm}.
\]
\end{answer}

\begin{answer}
\gksolhead{13}{解答题}
\gkanswer{答案 (1) $E_{0}=\dfrac{mgd(R_{1}+R_{2})}{qR_{2}}$；(2) $B=\dfrac{mv}{2qd}$；(3) $E'_{\min}=\dfrac{mg}{2q}$}
(1) 小球在极板间做匀速圆周运动，说明静电力与重力平衡，合力只由洛伦兹力提供。设极板间电压为 $U$，则
\[
  q\frac{U}{d}=mg
  \;\Longrightarrow\;
  U=\frac{mgd}{q}.
\]
电容器与 $R_{2}$ 并联，由串联分压
\[
  U=\frac{R_{2}}{R_{1}+R_{2}}E_{0},
\]
故
\[
  E_{0}=\frac{mgd(R_{1}+R_{2})}{qR_{2}}.
\]

(2) 小球从下板左边缘 $A$ 点进入，从上板右边缘离开，水平位移 $\sqrt{3}d$、竖直位移 $d$。设圆周运动半径为 $R$，圆心在 $A$ 点正上方，几何关系如图所示。

\gkfig{figures/s13-geometry}

由几何关系
\[
  (R-d)^{2}+(\sqrt{3}d)^{2}=R^{2},
\]
解得
\[
  R=2d .
\]
由洛伦兹力提供向心力
\[
  qvB=m\frac{v^{2}}{R},
\]
故
\[
  B=\frac{mv}{qR}=\frac{mv}{2qd}.
\]

(3) 由几何关系，小球离开电容器时速度方向与竖直方向的夹角为 $30^{\circ}$。小球离开电容器后只受重力和静电力，要做直线运动，合力必须沿速度方向，即静电力垂直于速度方向的分力必须与重力垂直于速度方向的分力等大反向。

当静电力沿速度方向的分力为零时，静电力最小，电场强度最小。此时
\[
  qE'=mg\sin30^{\circ},
\]
故
\[
  E'_{\min}=\frac{mg}{2q}.
\]
\end{answer}

\begin{answer}
\gksolhead{14}{解答题}
\gkanswer{答案 (1) $k=\sqrt{\dfrac{(1+\lambda)h}{(1-\lambda)H}}$；(2) $F_{0}=\dfrac{2(1-\lambda)(H-h)mg}{h-h_{0}}$}
\gkanswer{(3) $I=m\sqrt{\dfrac{2g(1-\lambda)(H-h)(H^{N+1}-h^{N+1})}{h(H^{N}-h^{N})}}$}
(1) 设篮球第一次碰地前速率为 $v_{\text{下}}$、碰后速率为 $v_{\text{上}}$。下落过程重力做正功、阻力做负功，由动能定理
\[
  (mg-\lambda mg)H=\frac{1}{2}mv_{\text{下}}^{2},
\]
反弹上升过程重力和阻力都做负功
\[
  -(mg+\lambda mg)h=0-\frac{1}{2}mv_{\text{上}}^{2}.
\]
解得
\[
  v_{\text{下}}=\sqrt{2(1-\lambda)gH},\qquad
  v_{\text{上}}=\sqrt{2(1+\lambda)gh},
\]
故碰后与碰前速率之比
\[
  k=\frac{v_{\text{上}}}{v_{\text{下}}}
  =\sqrt{\frac{(1+\lambda)h}{(1-\lambda)H}} .
\]

(2) 设第二次下落（从 $h$ 到地面）碰地前速率为 $v_{1}$、碰后速率为 $v_{2}$。下落过程中压力 $F$ 做的功等于 $F$--$y$ 图线与横轴所围面积
\[
  W_{F}=\frac{1}{2}F_{0}(h-h_{0}).
\]
由动能定理
\[
  \frac{1}{2}mv_{1}^{2}=(1-\lambda)mgh+W_{F}.
\]
碰后速率满足 $v_{2}=kv_{1}$，而反弹到 $h$ 需
\[
  \frac{1}{2}mv_{2}^{2}=(1+\lambda)mgh .
\]
联立得
\[
  (1+\lambda)mgh=k^{2}\left[(1-\lambda)mgh+W_{F}\right].
\]
代入 $k^{2}=\dfrac{(1+\lambda)h}{(1-\lambda)H}$，两边约去 $(1+\lambda)h$：
\[
  mgH=(1-\lambda)mgh+W_{F}
  \;\Longrightarrow\;
  W_{F}=(1-\lambda)mg(H-h).
\]
故
\[
  \frac{1}{2}F_{0}(h-h_{0})=(1-\lambda)mg(H-h),
\]
\[
  F_{0}=\frac{2(1-\lambda)(H-h)mg}{h-h_{0}} .
\]

(3) 每次拍击时间极短，重力的冲量可忽略，由动量定理，拍击瞬间给篮球竖直向下的速度
\[
  v=\frac{I}{m}.
\]
设第 $i$ 次拍击后篮球反弹的最大高度为 $h_{i}$（$h_{0}=h$，$h_{N}=H$）。从 $h_{i}$ 下落，碰地前速率满足
\[
  v_{i}^{2}=v^{2}+2(1-\lambda)gh_{i},
\]
碰后速率 $kv_{i}$，反弹高度由
\[
  (kv_{i})^{2}=2(1+\lambda)gh_{i+1}
\]
决定。代入 $k^{2}$ 得
\[
  h_{i+1}=\frac{h}{H}h_{i}+\frac{h}{H}\cdot\frac{v^{2}}{2(1-\lambda)g}.
\]
令 $a=\dfrac{h}{H}$，上式化为等比数列递推
\[
  h_{i+1}=a\,h_{i}+a\,\frac{v^{2}}{2(1-\lambda)g},
\]
迭代 $N$ 次并用 $h_{0}=h$：
\[
  h_{N}=a^{N}h+\frac{a(1-a^{N})}{1-a}\cdot\frac{v^{2}}{2(1-\lambda)g}.
\]
令 $h_{N}=H$，解得
\[
  v^{2}=\frac{2(1-\lambda)g(H-h)(H^{N+1}-h^{N+1})}{h(H^{N}-h^{N})},
\]
故
\[
  I=mv
  =m\sqrt{\frac{2g(1-\lambda)(H-h)(H^{N+1}-h^{N+1})}{h(H^{N}-h^{N})}} .
\]
\end{answer}

\begin{answer}
\gksolhead{15}{选考题（选修 3-3）}
\gkanswer{答案 (1) ABE；(2)(ⅰ) $p_{1}=1.0\times10^{5}$~Pa；(ⅱ) $F=1$~N}
(1) 分子热运动速率较小的气体从 $A$ 端流出，速率较大的从 $B$ 端流出；同种气体中分子平均速率较大的部分温度较高，故 $A$ 端为冷端、$B$ 端为热端，A 正确。由上述分析，$A$ 端流出的气体分子平均速率一定小于 $B$ 端流出的，B 正确。内能不仅与温度有关，还与气体分子数有关，题目未给出两端气体的分子数，不能判定 $A$ 端气体内能一定较大，C 错误。该装置中高压气体经喷嘴膨胀、分离冷热气，并非自发过程，不违背热力学第二定律；任何过程都满足能量守恒定律，故 D 错误、E 正确。

(2)(ⅰ) 活塞处于 $A$ 位置时，对活塞与 U 形金属丝整体受力分析，大气压力向上，内部气体压力与重力向下：
\[
  p_{0}S=p_{1}S+(m_{1}+m_{2})g,
\]
\[
  p_{1}=p_{0}-\frac{(m_{1}+m_{2})g}{S}
  =1.01\times10^{5}-\frac{0.10\times10}{10\times10^{-4}}~\si{\pascal}
  =1.0\times10^{5}~\si{\pascal}.
\]

(ⅱ) 活塞处于 $B$ 位置时，气体体积
\[
  V_{2}=V_{0}+Sh=9.9~\si{\litre}+10\times10^{-4}\times0.10\times10^{3}~\si{\litre}
  =10.0~\si{\litre}.
\]
环境温度不变，由玻意耳定律
\[
  p_{1}V_{0}=p_{2}V_{2},
\]
\[
  p_{2}=\frac{9.9}{10.0}\times1.0\times10^{5}~\si{\pascal}
  =9.9\times10^{4}~\si{\pascal}.
\]
此时对活塞与金属丝整体受力分析，液体对金属丝的拉力 $F$ 向下：
\[
  p_{0}S=p_{2}S+(m_{1}+m_{2})g+F,
\]
\[
  F=(p_{0}-p_{2})S-(m_{1}+m_{2})g
  =(1.01-0.99)\times10^{5}\times10\times10^{-4}~\si{\newton}-1.0~\si{\newton}
  =1~\si{\newton}.
\]
\end{answer}

\begin{answer}
\gksolhead{16}{选考题（选修 3-4）}
\gkanswer{答案 (1) ABD；(2)(ⅰ) $d=\dfrac{2\sqrt{15}}{5}$~mm；(ⅱ) $x=\dfrac{10\sqrt{3}-12}{15}$~mm}
(1) 木棒竖直方向做简谐运动，水平方向匀速运动。由 $F$--$x$ 图像知，浮力变化的“空间周期”为 \SI{0.4}{\metre}；振动周期 $T=1/f=\SI{1}{\second}$，故水平速度
\[
  v=\frac{0.4}{1}~\si{\metre\per\second}=\SI{0.4}{\metre\per\second}.
\]

A：$x$ 从 \SI{0.05}{\metre} 到 \SI{0.15}{\metre} 的过程中，浮力先减小后继续减小，木棒由深处向上运动，$x=\SI{0.1}{\metre}$ 处经过平衡位置，速度最大，故动能先增大后减小，A 正确。

B：$x$ 从 \SI{0.21}{\metre} 到 \SI{0.25}{\metre} 的过程中，浮力由最小值附近逐渐增大，木棒在平衡位置上方、正向下运动，重力大于浮力，合力向下且随浮力增大而减小，故加速度方向竖直向下、大小逐渐变小，B 正确。

C：$x=\SI{0.35}{\metre}$ 与 $x=\SI{0.45}{\metre}$ 关于 $x=\SI{0.4}{\metre}$ 对称，两处木棒竖直分速度等大反向，但水平分速度相同，合速度方向并不相反，C 错误。

D：设平衡时浸入水中的长度为 $h'$、振幅为 $A$，则
\[
  F_{1}=\rho gS(h'+A),\qquad F_{2}=\rho gS(h'-A),
\]
两式相减得
\[
  A=\frac{F_{1}-F_{2}}{2\rho gS},
\]
D 正确。

E：木棒是刚体，各质点随木棒一起整体运动，并非“振动在介质中传播”，不能视为机械波，E 错误。

(2)(ⅰ) 可视角度 $\theta=60^{\circ}$ 时，光线折射到空气后的最大折射角为 $30^{\circ}$。设介质中对应最大入射角为 $i$，光路如图所示。

\gkfig{figures/s16-screen}

由折射定律
\[
  \frac{\sin30^{\circ}}{\sin i}=n=2
  \;\Longrightarrow\;
  \sin i=\frac{1}{4},
\]
故
\[
  \tan i=\frac{1}{\sqrt{15}}.
\]
由几何关系，像素单元位于两屏障正中间，恰能从屏障顶端射出的光线满足
\[
  \tan i=\frac{l/2}{d},
\]
\[
  d=\frac{l/2}{\tan i}
  =\frac{0.4}{1/\sqrt{15}}~\si{\milli\metre}
  =\frac{2\sqrt{15}}{5}~\si{\milli\metre}.
\]

(ⅱ) 可视角度 $\theta'=180^{\circ}$ 时，最大折射角为 $90^{\circ}$，仍由折射定律
\[
  \frac{\sin90^{\circ}}{\sin i'}=2
  \;\Longrightarrow\;
  \sin i'=\frac{1}{2},
\]
故
\[
  \tan i'=\frac{1}{\sqrt{3}}.
\]
要使像素单元任意一点发出的光都能以 $90^{\circ}$ 折射出去，只需像素单元离某一屏障最远的一点满足临界条件。该点到屏障的水平距离为 $\dfrac{l}{2}+\dfrac{x}{2}$，则
\[
  \tan i'=\frac{l/2+x/2}{d'},
\]
\[
  \frac{l}{2}+\frac{x}{2}=d'\tan i'=\frac{1.0}{\sqrt{3}}~\si{\milli\metre}
  =\frac{\sqrt{3}}{3}~\si{\milli\metre},
\]
\[
  x=2\left(\frac{\sqrt{3}}{3}-\frac{l}{2}\right)
  =2\left(\frac{\sqrt{3}}{3}-0.4\right)~\si{\milli\metre}
  =\frac{10\sqrt{3}-12}{15}~\si{\milli\metre}.
\]
\end{answer}
